

\addcontentsline{toc}{section}{Введение}
\section*{Введение}
В различных сферах деятельности люди сталкиваются с необходимостью управления большим количеством данных, которые необходимо хранить, изменять и анализировать. Традиционные методы, такие как ведение бумажных архивов или использование простых электронных таблиц, в условиях больших объёмов часто приводят к долгому поиску нужных сведений, потере данных, ошибкам при обновлении записей и неструктурированным данным.

Решением этих проблем служат системы управления базами данных (СУБД) --- ключевой инструмент для структурирования и автоматизированной обработки данных. СУБД обеспечивает быстрый доступ к необходимым сведениям, ускоряет выполнение рутинных операций, таких как добавление или изменение, и, в отличие от локальных таблиц, позволяет нескольким пользователям одновременно работать с одной базой данных.

Перед программированием базы данных необходимо тщательно ее спроектировать, детально проанализировав предметную область, выделив ключевые сущности (таблицы) и определив атрибуты, значения которых будут в них храниться. Грамотно построенная структура таблиц позволяет минимизировать дублирование данных, гарантирует их целостность и безопасность, а также значительно повышает производительность всей информационной системы.

Современные СУБД предоставляют широкий спектр возможностей для управления данными и администрирования. Язык запросов SQL предоставляет возможность оперативно получать любые необходимые выборки по заданным критериям, а также модифицировать базу данных (добавлять, удалять, изменять или объединять записи). Настройка автоматического резервного копирования данных защищает данные от непредвиденных сбоев. Возможность предоставлять определенные права доступа тем или иным пользователям обеспечивает безопасную работу в системе.